% Options for packages loaded elsewhere
\PassOptionsToPackage{unicode}{hyperref}
\PassOptionsToPackage{hyphens}{url}
%
\documentclass[
]{article}
\usepackage{amsmath,amssymb}
\usepackage{iftex}
\ifPDFTeX
  \usepackage[T1]{fontenc}
  \usepackage[utf8]{inputenc}
  \usepackage{textcomp} % provide euro and other symbols
\else % if luatex or xetex
  \usepackage{unicode-math} % this also loads fontspec
  \defaultfontfeatures{Scale=MatchLowercase}
  \defaultfontfeatures[\rmfamily]{Ligatures=TeX,Scale=1}
\fi
\usepackage{lmodern}
\ifPDFTeX\else
  % xetex/luatex font selection
\fi
% Use upquote if available, for straight quotes in verbatim environments
\IfFileExists{upquote.sty}{\usepackage{upquote}}{}
\IfFileExists{microtype.sty}{% use microtype if available
  \usepackage[]{microtype}
  \UseMicrotypeSet[protrusion]{basicmath} % disable protrusion for tt fonts
}{}
\makeatletter
\@ifundefined{KOMAClassName}{% if non-KOMA class
  \IfFileExists{parskip.sty}{%
    \usepackage{parskip}
  }{% else
    \setlength{\parindent}{0pt}
    \setlength{\parskip}{6pt plus 2pt minus 1pt}}
}{% if KOMA class
  \KOMAoptions{parskip=half}}
\makeatother
\usepackage{xcolor}
\usepackage[margin=1in]{geometry}
\usepackage{graphicx}
\makeatletter
\def\maxwidth{\ifdim\Gin@nat@width>\linewidth\linewidth\else\Gin@nat@width\fi}
\def\maxheight{\ifdim\Gin@nat@height>\textheight\textheight\else\Gin@nat@height\fi}
\makeatother
% Scale images if necessary, so that they will not overflow the page
% margins by default, and it is still possible to overwrite the defaults
% using explicit options in \includegraphics[width, height, ...]{}
\setkeys{Gin}{width=\maxwidth,height=\maxheight,keepaspectratio}
% Set default figure placement to htbp
\makeatletter
\def\fps@figure{htbp}
\makeatother
\setlength{\emergencystretch}{3em} % prevent overfull lines
\providecommand{\tightlist}{%
  \setlength{\itemsep}{0pt}\setlength{\parskip}{0pt}}
\setcounter{secnumdepth}{-\maxdimen} % remove section numbering
\usepackage{graphicx}
\usepackage{float}
\usepackage{tikz}
\usepackage{xcolor}
\ifLuaTeX
  \usepackage{selnolig}  % disable illegal ligatures
\fi
\usepackage{bookmark}
\IfFileExists{xurl.sty}{\usepackage{xurl}}{} % add URL line breaks if available
\urlstyle{same}
\hypersetup{
  hidelinks,
  pdfcreator={LaTeX via pandoc}}

\author{}
\date{\vspace{-2.5em}}

\begin{document}

\section{Question}\label{question}

La gráfica muestra las exportaciones del sector minero de un país hacia
diferentes destinos durante el período 2008-2010, expresadas en millones
de dólares. Los principales productos exportados incluyen carbón y
petróleo, entre otros.

\includegraphics[width=0.85\textwidth,height=\textheight]{grafico_exportaciones_principal.png}

Otra representación que muestra toda la información de la gráfica
anterior es:

\textbf{A.}
\includegraphics[width=0.7\textwidth,height=\textheight]{opcion_a.png}

\textbf{B.}
\includegraphics[width=0.7\textwidth,height=\textheight]{opcion_b.png}

\textbf{C.}
\includegraphics[width=0.7\textwidth,height=\textheight]{opcion_c.png}

\textbf{D.}
\includegraphics[width=0.7\textwidth,height=\textheight]{opcion_d.png}

\subsection{Answerlist}\label{answerlist}

\begin{itemize}
\tightlist
\item
  La opción A representa mejor los datos porque muestra claramente los
  valores de 2010 con mayor detalle.
\item
  La opción C es correcta porque utiliza colores más llamativos que
  facilitan la interpretación visual.
\item
  La opción C muestra correctamente toda la información de la gráfica
  original, incluyendo los tres años (2008, 2009, 2010) y todos los
  países de destino.
\item
  La opción B es la más adecuada ya que permite visualizar la
  distribución porcentual entre países de destino.
\end{itemize}

\section{Solution}\label{solution}

Para resolver este problema, debemos analizar qué representación gráfica
muestra \textbf{toda la información} de la gráfica original de manera
equivalente.

\subsubsection{Análisis de la gráfica
original:}\label{anuxe1lisis-de-la-gruxe1fica-original}

La gráfica principal muestra: - \textbf{Datos de tres años}: 2008, 2009
y 2010 - \textbf{Cinco destinos}: Resto del mundo, China, Japón, Corea
del Sur, India - \textbf{Valores específicos} para cada combinación
año-destino - \textbf{Comparación temporal} entre los diferentes años

\subsubsection{Análisis de cada
opción:}\label{anuxe1lisis-de-cada-opciuxf3n}

\textbf{Opción A}: Muestra únicamente los datos de 2010 en formato de
barras horizontales. Aunque mantiene la información de todos los
destinos para ese año específico, \textbf{no incluye} los datos de 2008
y 2009.

\textbf{Opción B}: Presenta una gráfica circular con los datos de 2010,
mostrando la distribución porcentual entre destinos. \textbf{No incluye}
información temporal de los otros años.

\textbf{Opción C}: Muestra la evolución temporal de las exportaciones
mediante líneas, incluyendo \textbf{todos los años} (2008, 2009, 2010) y
\textbf{todos los destinos}. Esta representación permite visualizar
tanto los valores específicos como las tendencias temporales.

\textbf{Opción D}: Presenta únicamente los datos de 2009 en formato de
barras verticales. \textbf{No incluye} los datos de 2008 y 2010.

\subsubsection{Conclusión:}\label{conclusiuxf3n}

La \textbf{Opción C} es la única que muestra \textbf{toda la
información} de la gráfica original, incluyendo: - Los tres años de
datos (2008, 2009, 2010) - Los cinco destinos de exportación - Los
valores específicos para cada combinación - La posibilidad de comparar
tanto entre destinos como entre años

\subsection{Answerlist}\label{answerlist-1}

\begin{itemize}
\tightlist
\item
  Falso. Esta opción no considera todos los aspectos de la información
  presentada.
\item
  Falso. Esta opción no considera todos los aspectos de la información
  presentada.
\item
  Verdadero. Esta opción identifica correctamente la representación que
  muestra toda la información.
\item
  Falso. Esta opción no considera todos los aspectos de la información
  presentada.
\end{itemize}

\section{Meta-information}\label{meta-information}

exname: Exportaciones industriales interpretación representación extype:
schoice exsolution: 0010 exshuffle: TRUE exsection:
Estadística/Interpretación de datos

\end{document}
